\begin{abstract}

We study the relationship between topological phases of matter -- both topological insulators (TIs) and topologically ordered magnets -- and the geometry of the lattices on which they are constructed.

In the case of TIs, we consider the Chern number, the invariant used to classify Chern insulators, which is central to our understanding of non-interacting topological phases of matter. Such phases are generally highly robust to the presence of disorder, whereas the Chern number is dependent on translational symmetry to be well-defined. We consider several real-space formulations of the Chern number, showing how they are constructed, discussing their relative limitations and providing a novel physical interpretation for the real-space Chern number as a localised measurement of the Hall conductivity.

Next we study many-body topological order, focussing on Kitaev's honeycomb quantum spin liquid. We show that Kitaev's Hamiltonian can be constructed and exactly solved on almost any tri-coordinated lattice. Two complications arise in the non-crystalline case, namely the problems of the assigning bond labels and finding the ground state flux sector. We present solutions to both these problems and calculate the phase diagram on a general lattice. We show that a robust gapped quantum spin liquid ground state emerges with spontaneously broken time reversal symmetry, and discuss the stability of the system against bond-strength disorder.

Finally, we consider an extension of the Honeycomb model that includes the Heisenberg couplings present in candidate Kitaev materials. Here we consider the effect of geometric frustration, constructing the Kitaev-Heisenberg model on a non-bipartite lattice. In doing so, we frustrate antiferromagnetic order and find two valence bond solid phases that have no analogue on the honeycomb model. Triplon excitations are investigated in these phases, where we find that the interplay of an external magnetic field with the Kitaev couplings leads to a bosonic analogue of the quantum Hall effect, with a rich phase diagram of topological triplon bands. 

\end{abstract}